\documentclass[UTF8,12pt,fontset=windows]{ctexart}
\usepackage[paperwidth=240mm,paperheight=200mm,left=14mm,right=14mm,top=9mm,bottom=12mm]{geometry}
\usepackage{amsmath,amssymb,mathtools,bm,graphicx,xcolor}
\pagestyle{empty}
\setlength{\parindent}{0pt}
\setlength{\parskip}{4pt}
\setlength{\emergencystretch}{2em}
\xeCJKDeclareCharClass{CJK}{"2460 -> "24FF}
\xeCJKDeclareCharClass{CJK}{"2160 -> "2188}
\begin{document}
{\large\bfseries 2013.1}

设 $\cos x - 1 = x\sin\alpha(x)$，其中 $|\alpha(x)| < \frac{\pi}{2}$，则当 $x \to 0$ 时，$\alpha(x)$ 是（　）

（A）比 $x$ 高阶的无穷小量

（B）比 $x$ 低阶的无穷小量

（C）与 $x$ 同阶但不等价的无穷小量

（D）与 $x$ 等价的无穷小量

\vfill
\newpage
{\large\bfseries 2013.2}

设函数 $y = f(x)$ 由方程 $\cos(xy) + \ln y - x = 1$ 确定，则 $\lim_{n \to \infty}n\left[f\left(\frac{2}{n}\right) - 1\right] =$（　）

（A）2

（B）1

（C）$-1$

（D）$-2$

\vfill
\newpage
{\large\bfseries 2013.3}

设函数 $f(x) = \begin{cases} \sin x, & 0 \le x < \pi, \\ 2, & \pi \le x \le 2\pi, \end{cases}$ $F(x) = \int_0^x f(t)dt$，则（　）

（A）$x = \pi$ 是函数 $F(x)$ 的跳跃间断点

（B）$x = \pi$ 是函数 $F(x)$ 的可去间断点

（C）$F(x)$ 在 $x = \pi$ 处连续但不可导

（D）$F(x)$ 在 $x = \pi$ 处可导

\vfill
\newpage
{\large\bfseries 2013.4}

设函数 $f(x) = \begin{cases} \frac{1}{(x-1)^{\alpha-1}}, & 1 < x < e, \\ \frac{1}{x\ln^{\alpha+1}x}, & x \ge e, \end{cases}$ 若反常积分 $\int_1^{+\infty}f(x)dx$ 收敛，则（　）

（A）$\alpha < -2$

（B）$\alpha > 2$

（C）$-2 < \alpha < 0$

（D）$0 < \alpha < 2$

\vfill
\newpage
{\large\bfseries 2013.5}

设 $z = \frac{y}{x}f(xy)$，其中函数 $f$ 可微，则 $\frac{x}{y}\frac{\partial z}{\partial x} + \frac{\partial z}{\partial y} =$（　）

（A）$2yf'(xy)$

（B）$-2yf'(xy)$

（C）$\frac{2}{x}f(xy)$

（D）$-\frac{2}{x}f(xy)$

\vfill
\newpage
{\large\bfseries 2013.6}

设 $D_k$ 是圆域 $D = \{(x, y) \mid x^2 + y^2 \le 1\}$ 在第 $k$ 象限的部分，记 $I_k = \iint_{D_k}(y - x)dxdy$（$k = 1, 2, 3, 4$），则（　）

（A）$I_1 > 0$

（B）$I_2 > 0$

（C）$I_3 > 0$

（D）$I_4 > 0$

\vfill
\newpage
{\large\bfseries 2013.7}

设 $\boldsymbol{A}, \boldsymbol{B}, \boldsymbol{C}$ 均为 $n$ 阶矩阵。若 $\boldsymbol{AB} = \boldsymbol{C}$，且 $\boldsymbol{B}$ 可逆，则（　）

（A）矩阵 $\boldsymbol{C}$ 的行向量组与矩阵 $\boldsymbol{A}$ 的行向量组等价

（B）矩阵 $\boldsymbol{C}$ 的列向量组与矩阵 $\boldsymbol{A}$ 的列向量组等价

（C）矩阵 $\boldsymbol{C}$ 的行向量组与矩阵 $\boldsymbol{B}$ 的行向量组等价

（D）矩阵 $\boldsymbol{C}$ 的列向量组与矩阵 $\boldsymbol{B}$ 的列向量组等价

\vfill
\newpage
{\large\bfseries 2013.8}

矩阵 $\begin{pmatrix} 1 & a & 1 \\ a & b & a \\ 1 & a & 1 \end{pmatrix}$ 与 $\begin{pmatrix} 2 & 0 & 0 \\ 0 & b & 0 \\ 0 & 0 & 0 \end{pmatrix}$ 相似的充分必要条件为（　）

（A）$a = 0, b = 2$

（B）$a = 0, b$ 为任意常数

（C）$a = 2, b = 0$

（D）$a = 2, b$ 为任意常数

\vfill
\newpage
{\large\bfseries 2013.9}

$\lim_{x \to 0}\left[2 - \frac{\ln(1 + x)}{x}\right]^{\frac{1}{x}} =$ \underline{\hspace{20mm}}.

\vfill
\newpage
{\large\bfseries 2013.10}

设函数 $f(x) = \int_{-1}^x \sqrt{1 - e^t}dt$，则 $y = f(x)$ 的反函数 $x = f^{-1}(y)$ 在 $y = 0$ 处的导数 $\left.\frac{dx}{dy}\right|_{y = 0} =$ \underline{\hspace{20mm}}.

\vfill
\newpage
{\large\bfseries 2013.11}

设封闭曲线 $L$ 的极坐标方程为 $r = \cos 3\theta$（$-\frac{\pi}{6} \le \theta \le \frac{\pi}{6}$），则 $L$ 所围平面图形的面积是 \underline{\hspace{20mm}}.

\vfill
\newpage
{\large\bfseries 2013.12}

曲线 $\begin{cases} x = \arctan t, \\ y = \ln\sqrt{1 + t^2}, \end{cases}$ 上对应于 $t = 1$ 的点处的法线方程为 \underline{\hspace{20mm}}.

\vfill
\newpage
{\large\bfseries 2013.13}

已知 $y_1 = e^{3x} - xe^{2x}, y_2 = e^x - xe^{2x}, y_3 = -xe^{2x}$ 是某二阶常系数非齐次线性微分方程的3个解，则该方程满足条件 $y|_{x = 0} = 0, y'|_{x = 0} = 1$ 的解为 $y =$ \underline{\hspace{20mm}}.

\vfill
\newpage
{\large\bfseries 2013.14}

设 $\boldsymbol{A} = (a_{ij})$ 是3阶非零矩阵，$|\boldsymbol{A}|$ 为 $\boldsymbol{A}$ 的行列式，$A_{ij}$ 为 $a_{ij}$ 的代数余子式。若 $a_{ij} + A_{ij} = 0$（$i, j = 1, 2, 3$），则 $|\boldsymbol{A}| =$ \underline{\hspace{20mm}}.

\vfill
\newpage
{\large\bfseries 2013.15}

（本题满分10分）

当 $x \to 0$ 时，$1 - \cos x \cdot \cos 2x \cdot \cos 3x$ 与 $ax^n$ 为等价无穷小量，求 $n$ 与 $a$ 的值.

\vfill
\newpage
{\large\bfseries 2013.16}

（本题满分10分）

设 $D$ 是由曲线 $y = x^{\frac{1}{3}}$，直线 $x = a$（$a > 0$）及 $x$ 轴所围成的平面图形，$V_x, V_y$ 分别是 $D$ 绕 $x$ 轴、$y$ 轴旋转一周所得旋转体的体积。若 $V_y = 10V_x$，求 $a$ 的值.

\vfill
\newpage
{\large\bfseries 2013.17}

（本题满分10分）

设平面区域 $D$ 由直线 $x = 3y, y = 3x$ 及 $x + y = 8$ 围成，计算 $\iint_D x^2 dxdy$.

\vfill
\newpage
{\large\bfseries 2013.18}

（本题满分10分）

设奇函数 $f(x)$ 在 $[-1, 1]$ 上具有2阶导数，且 $f(1) = 1$. 证明：

（Ⅰ）存在 $\xi \in (0, 1)$，使得 $f'(\xi) = 1$；

（Ⅱ）存在 $\eta \in (-1, 1)$，使得 $f''(\eta) + f'(\eta) = 1$.

\vfill
\newpage
{\large\bfseries 2013.19}

（本题满分10分）

求曲线 $x^3 - xy + y^3 = 1$（$x \ge 0, y \ge 0$）上的点到坐标原点的最长距离和最短距离.

\vfill
\newpage
{\large\bfseries 2013.20}

（本题满分11分）

设函数 $f(x) = \ln x + \frac{1}{x}$，

（Ⅰ）求 $f(x)$ 的最小值；

（Ⅱ）设数列 $\{x_n\}$ 满足 $\ln x_n + \frac{1}{x_{n+1}} < 1$，证明 $\lim_{n \to \infty}x_n$ 存在，并求此极限.

\vfill
\newpage
{\large\bfseries 2013.21}

（本题满分11分）

设曲线 $L$ 的方程为 $y = \frac{1}{4}x^2 - \frac{1}{2}\ln x$（$1 \le x \le e$），

（Ⅰ）求 $L$ 的弧长；

（Ⅱ）设 $D$ 是由曲线 $L$，直线 $x = 1, x = e$ 及 $x$ 轴所围平面图形，求 $D$ 的形心的横坐标.

\vfill
\newpage
{\large\bfseries 2013.22}

（本题满分11分）

设 $\boldsymbol{A} = \begin{pmatrix} 1 & a \\ 1 & 0 \end{pmatrix}, \boldsymbol{B} = \begin{pmatrix} 0 & 1 \\ 1 & b \end{pmatrix}$。当 $a, b$ 为何值时，存在矩阵 $\boldsymbol{C}$ 使得 $\boldsymbol{AC} - \boldsymbol{CA} = \boldsymbol{B}$，并求所有矩阵 $\boldsymbol{C}$.

\vfill
\newpage
{\large\bfseries 2013.23}

（本题满分11分）

设二次型 $f(x_1, x_2, x_3) = 2(a_1x_1 + a_2x_2 + a_3x_3)^2 + (b_1x_1 + b_2x_2 + b_3x_3)^2$，记 $\boldsymbol{\alpha} = \begin{pmatrix} a_1 \\ a_2 \\ a_3 \end{pmatrix}, \boldsymbol{\beta} = \begin{pmatrix} b_1 \\ b_2 \\ b_3 \end{pmatrix}$.

（Ⅰ）证明二次型 $f$ 对应的矩阵为 $2\boldsymbol{\alpha}\boldsymbol{\alpha}^T + \boldsymbol{\beta}\boldsymbol{\beta}^T$；

（Ⅱ）若 $\boldsymbol{\alpha}, \boldsymbol{\beta}$ 正交且均为单位向量，证明 $f$ 在正交变换下的标准形为 $2y_1^2 + y_2^2$.

\vfill
\end{document}
